% =====================================================================
\chapter{Multimodalidad: imágenes y audio como tokens}
\label{tema:6}
% =====================================================================

Un Transformer opera sobre secuencias de vectores de dimensión $d$ y no depende de cómo se hayan obtenido. Para que procese imágenes o audio basta con una aplicación que convierta la señal en una secuencia de vectores, el análogo de la tabla de \emph{embeddings} para el texto. Este tema estudia las dos construcciones más citadas: el Vision Transformer, que divide la imagen en parches y proyecta cada uno linealmente, y los codificadores de audio, que parten de un espectrograma o de un códec neuronal que produce tokens discretos. Se sigue a Dosovitskiy \emph{et al.}~\cite[§3.1]{vit}, la §11.8 de Zhang \emph{et al.}~\cite{d2l}, a Radford \emph{et al.}~\cite[§2.1 y §2.2]{whisper} y a Zeghidour \emph{et al.}~\cite{soundstream}, junto con el informe técnico de Gemini~1.5~\cite[§3.1]{gemini15}.

\section{El Vision Transformer}

\begin{definicion}[Incrustación por parches]\label{def:vit}
Sea una imagen $X\in\R^{H_{\mathrm{img}}\times W_{\mathrm{img}}\times C_{\mathrm{in}}}$ y sea $P$ un divisor de $H_{\mathrm{img}}$ y de $W_{\mathrm{img}}$. Se divide la imagen en $N_p=H_{\mathrm{img}}W_{\mathrm{img}}/P^{2}$ parches de $P\times P$ píxeles, se aplana cada uno en un vector $x_p^{(k)}\in\R^{1\times P^{2}C_{\mathrm{in}}}$ y se define
\[
z_0=\bigl[x_{\mathrm{clase}};\ x_p^{(1)}E;\ x_p^{(2)}E;\ \dots;\ x_p^{(N_p)}E\bigr]+E_{\mathrm{pos}},
\]
con $E\in\R^{P^{2}C_{\mathrm{in}}\times d}$, un vector aprendido $x_{\mathrm{clase}}\in\R^{1\times d}$ y posiciones aprendidas $E_{\mathrm{pos}}\in\R^{(N_p+1)\times d}$~\cite[§3.1, ec.~1]{vit}. La secuencia $z_0$, de $N_p+1$ filas, entra en un Transformer codificador.
\end{definicion}

\begin{proposicion}[Tokens y coste de un ViT]\label{prop:vit}
Con la notación de la \cref{def:vit}:
\begin{enumerate}
  \item la proyección $E$ tiene $P^{2}C_{\mathrm{in}}d$ parámetros, independientemente del tamaño de la imagen;
  \item si el lado de la imagen se multiplica por $\kappa$ con $P$ fijo, el número de parches se multiplica por $\kappa^{2}$, el coste de las capas lineales de cada bloque por unos $\kappa^{2}$ y el de los productos de la atención por unos $\kappa^{4}$.
\end{enumerate}
\end{proposicion}

\begin{table}[ht]
\centering\small
\caption{Número de parches y dimensión de cada parche aplanado para imágenes RGB ($C_{\mathrm{in}}=3$).}\label{tab:vit}
\begin{tabular}{@{}cccc@{}}
\toprule
Imagen & Parche & $N_p$ & $P^{2}C_{\mathrm{in}}$ \\
\midrule
$224\times 224$ & $16\times 16$ & $14^{2}=196$ & 768 \\
$384\times 384$ & $16\times 16$ & $24^{2}=576$ & 768 \\
$224\times 224$ & $32\times 32$ & $7^{2}=49$ & 3072 \\
\bottomrule
\end{tabular}
\end{table}

El modelo ViT-Base tiene 12 capas, $d=768$, un MLP de dimensión 3072 y 12 cabezas, con unos $8{,}6\cdot 10^{7}$ parámetros~\cite[tabla~1]{vit}. Con parches de 16 y una imagen de $224\times 224$, la secuencia tiene $196+1=197$ tokens y $E$ tiene $768\cdot 768=589\,824$ parámetros.

\begin{observacion}[Parches y convoluciones]
Aplicar $E$ a cada parche es lo mismo que aplicar una convolución con $d$ filtros de tamaño $P\times P$ y paso $P$: cada filtro es una columna de $E$ reordenada como un tensor $P\times P\times C_{\mathrm{in}}$, y el paso igual al tamaño hace que los parches no se solapen. Es la forma en que suele implementarse.
\end{observacion}

\section{Audio}

El audio es una señal unidimensional muestreada a miles de muestras por segundo: a $16$~kHz, treinta segundos son $480\,000$ muestras, demasiadas para usarlas como tokens. Hay dos formas habituales de reducirlas.

\addtocounter{definicion}{1}%

La segunda vía convierte el audio en enteros. Un \emph{códec neuronal} como SoundStream~\cite{soundstream} o EnCodec~\cite{encodec} consta de un codificador convolucional que produce un vector por cada fotograma de unos milisegundos, un cuantizador que sustituye ese vector por índices de diccionarios finitos y un decodificador que reconstruye la onda. SoundStream usa un cuantizador vectorial residual: el primer diccionario aproxima el vector, el segundo cuantiza el error del primero, y así sucesivamente.

\begin{proposicion}[Tasa de un cuantizador vectorial residual]\label{prop:rvq}
Un códec que produce $\nu$ fotogramas por segundo y cuantiza cada uno con $N_q$ diccionarios de $K$ entradas genera $\nu N_q$ tokens por segundo y una tasa de $\nu N_q\log_2K$ bits por segundo.
\end{proposicion}

Con valores ilustrativos, $\nu=50$~Hz, $N_q=8$ y $K=1024$, se obtienen 400 tokens por segundo y $4$~kbit/s, dentro del intervalo de 3 a 18~kbit/s que cubre SoundStream~\cite{soundstream}. Con estos tokens, el audio se modela con el mismo objetivo de predicción del siguiente token del módulo~1.

\section{Fusión con el modelo de lenguaje}

Hay dos formas de unir los vectores de otra modalidad con los de texto. En la \emph{fusión con adaptador}, un codificador separado (un ViT, por ejemplo) produce vectores que una proyección lineal o un pequeño MLP lleva a la dimensión $d$ del modelo de lenguaje, y se insertan en la secuencia de texto. En la \emph{fusión nativa}, el modelo se entrena desde el principio con secuencias en las que se intercalan tokens de varias modalidades.

\begin{observacion}[Lo que es público]
El informe de Gemini~1.5 afirma que el modelo es multimodal de forma nativa y admite datos de distintas modalidades intercalados en una misma secuencia: texto, imágenes, audio y vídeo, con ejemplos de 107 horas de audio y 10,5 horas de vídeo a un fotograma por segundo~\cite[§3.1]{gemini15}. No describe sus codificadores. Claude y las versiones de Grok posteriores a Grok-1 aceptan imágenes según la documentación pública de sus productos, pero sus codificadores tampoco son públicos; Grok-1, el único con pesos publicados, es un modelo de texto~\cite{grok1}. Lo que sí puede afirmarse en general es que, sea cual sea el codificador, las demás modalidades acaban como vectores en el mismo flujo residual que el texto.
\end{observacion}

